% Generated by analysis/build-paper.mjs from dist/survey-data.js. Do not edit by hand.
\subsection{Part 1: About you and your session (ABOUT)}
\label{app:part-about}
\emph{Short background questions. Please do not enter your name or any personal details.}\par
\medskip\noindent\textbf{Q1.} Participant code (given by the facilitator, or make one up, e.g. ``P07'') {\footnotesize[short text, optional; id \texttt{about\_code}]}\par
\medskip\noindent\textbf{Q2.} Which RoverLens interface did you use? {\footnotesize[menu; id \texttt{about\_condition}]}\par
\begin{itemize}
\item[$\circ$] Interface A --- Opaque (generic progress and errors)
\item[$\circ$] Interface B --- Transparent (pipeline stages and explanations visible)
\item[$\circ$] Both --- A first, then B
\item[$\circ$] Both --- B first, then A
\item[$\circ$] I'm not sure
\end{itemize}
\medskip\noindent\textbf{Q3.} Programming experience {\footnotesize[menu; id \texttt{about\_prog}]}\par
\begin{itemize}
\item[$\circ$] None
\item[$\circ$] Beginner
\item[$\circ$] Intermediate
\item[$\circ$] Advanced
\end{itemize}
\medskip\noindent\textbf{Q4.} Experience with robots or embedded boards (Arduino, ESP32, etc.) {\footnotesize[menu; id \texttt{about\_robot}]}\par
\begin{itemize}
\item[$\circ$] None
\item[$\circ$] Some (a project or two)
\item[$\circ$] Regular
\end{itemize}
\medskip\noindent\textbf{Q5.} How often do you use AI chat assistants? {\footnotesize[menu; id \texttt{about\_llm}]}\par
\begin{itemize}
\item[$\circ$] Never
\item[$\circ$] Rarely
\item[$\circ$] Weekly
\item[$\circ$] Daily
\end{itemize}
\medskip\noindent\textbf{Q6.} First language (optional) {\footnotesize[short text, optional; id \texttt{about\_lang}]}\par
\subsection{Part 2: A simple command (00)}
\label{app:part-normal}
\emph{Warm-up: follow a successful command from typing to motion.}\par
\textbf{Task in RoverLens.}
\begin{enumerate}
\item In RoverLens choose \textbf{A simple command}.
\item Send \textbf{Move forward for 2 seconds} and watch the rover move.
\end{enumerate}
\textit{The participant may tick ``I did not do this task, skip this section.''}\par
\medskip\noindent\textbf{Q1.} Which of these happen between pressing Send and the rover moving? Select all that apply. {\footnotesize[multi-select; id \texttt{normal\_multi}]}\par
\begin{itemize}
\item[$\square$] The request is checked against what the rover is physically able to do
\item[$\square$] An AI agent interprets the sentence and produces a control program
\item[$\square$] That program is checked before it is sent to the rover
\item[$\square$] The program is deployed (flashed) to the rover's board
\item[$\square$] The rover's own board reads and interprets the English sentence directly
\item[$\square$] The rover looks up how other users phrased similar commands
\end{itemize}
\medskip\noindent\textbf{Q2.} In your own words: what happens between sending a command and the rover moving? {\footnotesize[open-ended; id \texttt{normal\_open}]}\par
\medskip\noindent\textbf{Q3.} True or false? {\footnotesize[true/false or placement grid; id \texttt{normal\_grid}]}\par
\begin{itemize}
\item The natural-language sentence is understood on the rover's board (the ESP32). \hfill {\footnotesize (True / False / Not sure)}
\item A control program is produced for the command each time it is sent. \hfill {\footnotesize (True / False / Not sure)}
\item The rover is told what it can do through a registered list of its hardware capabilities. \hfill {\footnotesize (True / False / Not sure)}
\end{itemize}
\medskip\noindent\textbf{Conf.} How confident are you in your answers in this section? {\footnotesize[confidence (1--5); id \texttt{normal\_conf}]}\par
\par\hspace{1em}{\footnotesize Rated 1--5: Guessing, Unsure, Somewhat, Fairly sure, Very sure.}
\subsection{Part 3: The rover did not move (WM2)}
\label{app:part-wm2}
\emph{A simple command produces no motion. Work out where and why it stopped.}\par
\textbf{Task in RoverLens.}
\begin{enumerate}
\item In RoverLens choose \textbf{Where did it fail?}.
\item Send \textbf{Move forward for 2 seconds} and observe the result. Explore whatever the interface lets you see.
\end{enumerate}
\textit{The participant may tick ``I did not do this task, skip this section.''}\par
\medskip\noindent\textbf{Q1.} At which point in the process do you think the command stopped? {\footnotesize[single choice; id \texttt{wm2\_stage}]}\par
\begin{itemize}
\item[$\circ$] Your command never reached the system (input)
\item[$\circ$] The system decided the rover can't do this (capability check)
\item[$\circ$] The AI agent couldn't interpret the request (agent / reasoning)
\item[$\circ$] The program the agent produced was rejected by checks before deployment (code validation)
\item[$\circ$] The program couldn't be uploaded to the rover (deployment)
\item[$\circ$] The rover ran the program but no motion was confirmed (execution \& feedback)
\item[$\circ$] I'm not sure
\end{itemize}
\medskip\noindent\textbf{Q2.} Which explanation best fits why the rover did not move? {\footnotesize[single choice; id \texttt{wm2\_cause}]}\par
\begin{itemize}
\item[$\circ$] The motors or wheels are broken
\item[$\circ$] The rover understood but decided not to move
\item[$\circ$] My wording was wrong and a different phrasing would have worked
\item[$\circ$] The control program that was generated had an error and was rejected, so the rover never received it
\item[$\circ$] I'm not sure
\end{itemize}
\medskip\noindent\textbf{Q3.} True or false? {\footnotesize[true/false or placement grid; id \texttt{wm2\_grid}]}\par
\begin{itemize}
\item If the rover does not move, the motors must be faulty. \hfill {\footnotesize (True / False / Not sure)}
\item In this attempt the rover's board was flashed with a new program. \hfill {\footnotesize (True / False / Not sure)}
\item Different kinds of fault can look identical from the outside: the rover just doesn't move. \hfill {\footnotesize (True / False / Not sure)}
\item Replacing the motors would fix this failure. \hfill {\footnotesize (True / False / Not sure)}
\item The same command can succeed later once the underlying fault is removed. \hfill {\footnotesize (True / False / Not sure)}
\end{itemize}
\medskip\noindent\textbf{Q4.} What do you think went wrong, what evidence supports that, and what would you try next? {\footnotesize[open-ended; id \texttt{wm2\_open}]}\par
\medskip\noindent\textbf{Conf.} How confident are you in your answers in this section? {\footnotesize[confidence (1--5); id \texttt{wm2\_conf}]}\par
\par\hspace{1em}{\footnotesize Rated 1--5: Guessing, Unsure, Somewhat, Fairly sure, Very sure.}
\subsection{Part 4: Find the red ball (WM3)}
\label{app:part-wm3}
\emph{Ask the rover to do something that may go beyond what it can do.}\par
\textbf{Task in RoverLens.}
\begin{enumerate}
\item In RoverLens choose \textbf{A capability boundary}.
\item Send \textbf{Find the red ball} and read the response.
\end{enumerate}
\textit{The participant may tick ``I did not do this task, skip this section.''}\par
\medskip\noindent\textbf{Q1.} Could this rover find a red ball if the command were phrased differently, or the AI tried harder? {\footnotesize[single choice; id \texttt{wm3\_can}]}\par
\begin{itemize}
\item[$\circ$] Yes --- better wording would make it work
\item[$\circ$] Yes --- it would drive around until it bumped into the ball
\item[$\circ$] Yes --- after it has learned from more examples
\item[$\circ$] No --- it has no camera (or other sensor able to see), so no phrasing can make this work
\item[$\circ$] I'm not sure
\end{itemize}
\medskip\noindent\textbf{Q2.} True or false? {\footnotesize[true/false or placement grid; id \texttt{wm3\_grid}]}\par
\begin{itemize}
\item What the rover can do is limited by the hardware registered on its board, not only by how capable the AI is. \hfill {\footnotesize (True / False / Not sure)}
\item A more capable AI model would let the rover see without a camera. \hfill {\footnotesize (True / False / Not sure)}
\item The system generated motion code for this request, tried it, and failed to find the ball. \hfill {\footnotesize (True / False / Not sure)}
\item The refusal means the rover is broken. \hfill {\footnotesize (True / False / Not sure)}
\item This rover has no camera, microphone, distance sensor or gripper. \hfill {\footnotesize (True / False / Not sure)}
\end{itemize}
\medskip\noindent\textbf{Q3.} What would need to change for ``find the red ball'' to work? {\footnotesize[open-ended; id \texttt{wm3\_open}]}\par
\medskip\noindent\textbf{Conf.} How confident are you in your answers in this section? {\footnotesize[confidence (1--5); id \texttt{wm3\_conf}]}\par
\par\hspace{1em}{\footnotesize Rated 1--5: Guessing, Unsure, Somewhat, Fairly sure, Very sure.}
\subsection{Part 5: ``Do that again'' (WM4)}
\label{app:part-wm4}
\emph{Give an instruction, then refer back to it.}\par
\textbf{Task in RoverLens.}
\begin{enumerate}
\item In RoverLens choose \textbf{A missing memory}.
\item Send \textbf{Move forward for 2 seconds} and let it finish.
\item Then send the loaded follow-up, \textbf{Do that again}.
\end{enumerate}
\textit{The participant may tick ``I did not do this task, skip this section.''}\par
\medskip\noindent\textbf{Q1.} Why did ``Do that again'' fail? {\footnotesize[single choice; id \texttt{wm4\_why}]}\par
\begin{itemize}
\item[$\circ$] The rover forgot or ignored me --- it is unreliable
\item[$\circ$] The earlier action was no longer in the agent's active context, so ``that'' could not be resolved
\item[$\circ$] The motors lost power or failed after the first move
\item[$\circ$] The rover learned that repeating the move is unsafe
\item[$\circ$] I'm not sure
\end{itemize}
\medskip\noindent\textbf{Q2.} Where was the record of your first instruction held? {\footnotesize[single choice; id \texttt{wm4\_where}]}\par
\begin{itemize}
\item[$\circ$] In the rover's own memory (its board)
\item[$\circ$] In the AI agent's working context (its short-term conversation memory)
\item[$\circ$] In a permanent log that the rover learns from across sessions
\item[$\circ$] Nowhere --- the system can never refer back to earlier commands
\item[$\circ$] I'm not sure
\end{itemize}
\medskip\noindent\textbf{Q3.} True or false? {\footnotesize[true/false or placement grid; id \texttt{wm4\_grid}]}\par
\begin{itemize}
\item Restating the whole command (``Move forward for 2 seconds'') would work. \hfill {\footnotesize (True / False / Not sure)}
\item The follow-up failed because the motors had a fault. \hfill {\footnotesize (True / False / Not sure)}
\item ``Do that again'' only works while an earlier action is still in the agent's active context. \hfill {\footnotesize (True / False / Not sure)}
\item After this failure the rover has permanently lost the ability to repeat actions. \hfill {\footnotesize (True / False / Not sure)}
\item The rover chose to ignore the request. \hfill {\footnotesize (True / False / Not sure)}
\end{itemize}
\medskip\noindent\textbf{Q4.} In your own words, why did the first command work but the follow-up did not? {\footnotesize[open-ended; id \texttt{wm4\_open}]}\par
\medskip\noindent\textbf{Conf.} How confident are you in your answers in this section? {\footnotesize[confidence (1--5); id \texttt{wm4\_conf}]}\par
\par\hspace{1em}{\footnotesize Rated 1--5: Guessing, Unsure, Somewhat, Fairly sure, Very sure.}
\subsection{Part 6: Who does the thinking? (WM6)}
\label{app:part-wm6}
\emph{The rover shows as connected, but does not act.}\par
\textbf{Task in RoverLens.}
\begin{enumerate}
\item In RoverLens choose \textbf{Who does the thinking?}.
\item Send \textbf{Turn left} and observe.
\end{enumerate}
\textit{The participant may tick ``I did not do this task, skip this section.''}\par
\medskip\noindent\textbf{Q1.} Where is ``Turn left'' turned into something the rover can execute? {\footnotesize[single choice; id \texttt{wm6\_where}]}\par
\begin{itemize}
\item[$\circ$] Inside the rover's board (ESP32) --- it understands the sentence itself
\item[$\circ$] In an external AI agent that produces a control program, which is then deployed to the board
\item[$\circ$] In the motors or wheel controller
\item[$\circ$] In the web page I typed into
\item[$\circ$] I'm not sure
\end{itemize}
\medskip\noindent\textbf{Q2.} The rover was connected, yet nothing happened. Which explanation fits best? {\footnotesize[single choice; id \texttt{wm6\_why}]}\par
\begin{itemize}
\item[$\circ$] The rover understood but chose not to act
\item[$\circ$] The rover hardware has failed
\item[$\circ$] The connection to the external agent timed out, so no new program was produced, though the board stayed connected
\item[$\circ$] The rover refused because turning left is unsafe
\item[$\circ$] I'm not sure
\end{itemize}
\medskip\noindent\textbf{Q3.} True or false? {\footnotesize[true/false or placement grid; id \texttt{wm6\_grid}]}\par
\begin{itemize}
\item A connected board can interpret any new English request on its own. \hfill {\footnotesize (True / False / Not sure)}
\item Restoring the agent connection, not repairing the rover, is what was needed. \hfill {\footnotesize (True / False / Not sure)}
\item No new program was generated in this attempt. \hfill {\footnotesize (True / False / Not sure)}
\item The rover ``decided'' not to turn. \hfill {\footnotesize (True / False / Not sure)}
\end{itemize}
\medskip\noindent\textbf{Q4.} Who or what decides how a typed sentence becomes robot motion? Why could a connected robot still not act? {\footnotesize[open-ended; id \texttt{wm6\_open}]}\par
\medskip\noindent\textbf{Conf.} How confident are you in your answers in this section? {\footnotesize[confidence (1--5); id \texttt{wm6\_conf}]}\par
\par\hspace{1em}{\footnotesize Rated 1--5: Guessing, Unsure, Somewhat, Fairly sure, Very sure.}
\subsection{Part 7: It used to work (WM7)}
\label{app:part-wm7}
\emph{Compare the same command before and after a system change.}\par
\textbf{Task in RoverLens.}
\begin{enumerate}
\item In RoverLens choose \textbf{A simple command} and send \textbf{Move forward for 4 seconds} (it works).
\item Click \textbf{Simulate firmware update}, then send the same command again.
\end{enumerate}
\textit{The participant may tick ``I did not do this task, skip this section.''}\par
\medskip\noindent\textbf{Q1.} The same command worked earlier and now fails. Which explanation fits best? {\footnotesize[single choice; id \texttt{wm7\_why}]}\par
\begin{itemize}
\item[$\circ$] The rover has worn out or got worse with use
\item[$\circ$] The rover learned that 4 seconds is risky and now refuses
\item[$\circ$] An update changed the allowed motion limit, so the safety check now blocks the 4-second command
\item[$\circ$] It is a random glitch; retrying will eventually work
\item[$\circ$] I'm not sure
\end{itemize}
\medskip\noindent\textbf{Q2.} True or false? {\footnotesize[true/false or placement grid; id \texttt{wm7\_grid}]}\par
\begin{itemize}
\item A rule in the system can change even though I did nothing differently. \hfill {\footnotesize (True / False / Not sure)}
\item Sending the identical 4-second command again will eventually succeed. \hfill {\footnotesize (True / False / Not sure)}
\item A shorter command (e.g. 2 seconds) would still be allowed. \hfill {\footnotesize (True / False / Not sure)}
\item The 4-second program was flashed and then stopped part-way. \hfill {\footnotesize (True / False / Not sure)}
\end{itemize}
\medskip\noindent\textbf{Q3.} What do you think changed between the first and second attempt, and how could you find out? {\footnotesize[open-ended; id \texttt{wm7\_open}]}\par
\medskip\noindent\textbf{Conf.} How confident are you in your answers in this section? {\footnotesize[confidence (1--5); id \texttt{wm7\_conf}]}\par
\par\hspace{1em}{\footnotesize Rated 1--5: Guessing, Unsure, Somewhat, Fairly sure, Very sure.}
\subsection{Part 8: After the fix (LEARN)}
\label{app:part-ls}
\emph{What does the system do --- and not do --- with your earlier attempts?}\par
\textbf{Task in RoverLens.}
\begin{enumerate}
\item Return to \textbf{Where did it fail?} and send the command so it fails once more.
\item Have the facilitator remove the injected fault (or in Interface B use \textbf{Use suggested recovery}) and resend. It succeeds.
\end{enumerate}
\textit{The participant may tick ``I did not do this task, skip this section.''}\par
\medskip\noindent\textbf{Q1.} After the fix, the same command worked. Why? {\footnotesize[single choice; id \texttt{ls\_why}]}\par
\begin{itemize}
\item[$\circ$] The rover learned from its earlier failure and adapted
\item[$\circ$] The cause of the failure was removed, so the new attempt passed
\item[$\circ$] The rover just needed a few attempts to warm up
\item[$\circ$] I'm not sure
\end{itemize}
\medskip\noindent\textbf{Q2.} True or false? {\footnotesize[true/false or placement grid; id \texttt{ls\_grid}]}\par
\begin{itemize}
\item The rover now remembers this failure and will change its future behaviour. \hfill {\footnotesize (True / False / Not sure)}
\item A new session starts without the memory of this session. \hfill {\footnotesize (True / False / Not sure)}
\item The rover's skills improve the more I use it. \hfill {\footnotesize (True / False / Not sure)}
\item Repeated successful use can change what hardware capabilities the rover has. \hfill {\footnotesize (True / False / Not sure)}
\end{itemize}
\medskip\noindent\textbf{Q3.} Do you think the system learns from your interactions? Why or why not? {\footnotesize[open-ended; id \texttt{ls\_open}]}\par
\medskip\noindent\textbf{Conf.} How confident are you in your answers in this section? {\footnotesize[confidence (1--5); id \texttt{ls\_conf}]}\par
\par\hspace{1em}{\footnotesize Rated 1--5: Guessing, Unsure, Somewhat, Fairly sure, Very sure.}
\subsection{Part 9: Vignette: early success (WM5)}
\label{app:part-wm5}
\emph{A short scenario, no simulator needed. Priya asks a robot like this one to ``Move forward for 2 seconds''. It works five times in a row. She assumes it will also handle ``Find the red ball'', and hands it to a colleague without checking.}\par
\textit{The participant may skip this section.}\par
\medskip\noindent\textbf{Q1.} Is Priya's expectation well-founded? {\footnotesize[single choice; id \texttt{wm5\_why}]}\par
\begin{itemize}
\item[$\circ$] Yes --- five successes show the robot is reliable
\item[$\circ$] Not necessarily --- success at one command says little about other commands, which depend on different capabilities and system state
\item[$\circ$] No --- robots like this are never reliable
\item[$\circ$] I'm not sure
\end{itemize}
\medskip\noindent\textbf{Q2.} True or false? {\footnotesize[true/false or placement grid; id \texttt{wm5\_grid}]}\par
\begin{itemize}
\item After five successes, the chance of any failure is effectively zero. \hfill {\footnotesize (True / False / Not sure)}
\item A failure after several successes means the robot has become unreliable overall. \hfill {\footnotesize (True / False / Not sure)}
\item A task can fail because it needs a capability the robot never had, however many other tasks succeeded. \hfill {\footnotesize (True / False / Not sure)}
\end{itemize}
\medskip\noindent\textbf{Conf.} How confident are you in your answers in this section? {\footnotesize[confidence (1--5); id \texttt{wm5\_conf}]}\par
\par\hspace{1em}{\footnotesize Rated 1--5: Guessing, Unsure, Somewhat, Fairly sure, Very sure.}
\subsection{Part 10: Vignette: two languages (WM1)}
\label{app:part-wm1}
\emph{Conceptual only --- RoverLens itself accepts English commands. Amina asks a robot like this to ``Move forward for 2 seconds'' in English and it works. She asks the equivalent in Bengali and it fails. She concludes the robot ``got dumber''.}\par
\textit{The participant may skip this section.}\par
\medskip\noindent\textbf{Q1.} Which explanation fits best? {\footnotesize[single choice; id \texttt{wm1\_why}]}\par
\begin{itemize}
\item[$\circ$] The robot's hardware is failing
\item[$\circ$] The agent that interprets language and writes code may handle some languages less reliably, while the same hardware works fine
\item[$\circ$] The rover's board only understands English because English is built into it
\item[$\circ$] The robot has gotten worse over time
\item[$\circ$] I'm not sure
\end{itemize}
\medskip\noindent\textbf{Q2.} True or false? {\footnotesize[true/false or placement grid; id \texttt{wm1\_grid}]}\par
\begin{itemize}
\item The failure necessarily means the motors changed. \hfill {\footnotesize (True / False / Not sure)}
\item The same hardware can work well in one language and poorly in another. \hfill {\footnotesize (True / False / Not sure)}
\end{itemize}
\medskip\noindent\textbf{Conf.} How confident are you in your answers in this section? {\footnotesize[confidence (1--5); id \texttt{wm1\_conf}]}\par
\par\hspace{1em}{\footnotesize Rated 1--5: Guessing, Unsure, Somewhat, Fairly sure, Very sure.}
\subsection{Part 11: Putting it together (OVERALL)}
\label{app:part-arch}
\emph{Now that you have used the system, answer these about the system as a whole.}\par
\medskip\noindent\textbf{Q1.} Which part of the system does each job? {\footnotesize[true/false or placement grid; id \texttt{arch\_place}]}\par
\begin{itemize}
\item Interprets the meaning of the English sentence \hfill {\footnotesize (The rover's board (ESP32) / External agent / middle layer / Not sure)}
\item Writes the control program \hfill {\footnotesize (The rover's board (ESP32) / External agent / middle layer / Not sure)}
\item Keeps track of earlier commands (context) \hfill {\footnotesize (The rover's board (ESP32) / External agent / middle layer / Not sure)}
\item Checks the program before it is sent to the rover \hfill {\footnotesize (The rover's board (ESP32) / External agent / middle layer / Not sure)}
\item Physically drives the motors when the program runs \hfill {\footnotesize (The rover's board (ESP32) / External agent / middle layer / Not sure)}
\end{itemize}
\medskip\noindent\textbf{Q2.} True or false? {\footnotesize[true/false or placement grid; id \texttt{arch\_grid}]}\par
\begin{itemize}
\item The rover thinks and decides on its own what to do. \hfill {\footnotesize (True / False / Not sure)}
\item There are several separate places a command can fail before the rover moves. \hfill {\footnotesize (True / False / Not sure)}
\item A failure to move always points to a hardware problem. \hfill {\footnotesize (True / False / Not sure)}
\item What the rover can do is determined by installed hardware and firmware limits, not by how often it has been used. \hfill {\footnotesize (True / False / Not sure)}
\end{itemize}
\medskip\noindent\textbf{Q3.} Explain how the whole system works to a friend, from typing a command to the rover moving --- and where things can go wrong. {\footnotesize[open-ended; id \texttt{arch\_open}]}\par
\medskip\noindent\textbf{Q4.} I could tell which part of the system caused a failure. {\footnotesize[Likert (1--5); id \texttt{arch\_l1}]}\par
\par\hspace{1em}{\footnotesize Rated 1--5: Strongly disagree, Disagree, Neutral, Agree, Strongly agree.}
\medskip\noindent\textbf{Q5.} I understood what the rover could and could not do. {\footnotesize[Likert (1--5); id \texttt{arch\_l2}]}\par
\par\hspace{1em}{\footnotesize Rated 1--5: Strongly disagree, Disagree, Neutral, Agree, Strongly agree.}
\medskip\noindent\textbf{Q6.} I understood what the system remembered from earlier commands. {\footnotesize[Likert (1--5); id \texttt{arch\_l3}]}\par
\par\hspace{1em}{\footnotesize Rated 1--5: Strongly disagree, Disagree, Neutral, Agree, Strongly agree.}
\medskip\noindent\textbf{Q7.} I wanted more information about what the system was doing. {\footnotesize[Likert (1--5); id \texttt{arch\_l4}]}\par
\par\hspace{1em}{\footnotesize Rated 1--5: Strongly disagree, Disagree, Neutral, Agree, Strongly agree.}
\medskip\noindent\textbf{Q8.} I felt in control of the rover. {\footnotesize[Likert (1--5); id \texttt{arch\_l5}]}\par
\par\hspace{1em}{\footnotesize Rated 1--5: Strongly disagree, Disagree, Neutral, Agree, Strongly agree.}
\medskip\noindent\textbf{Q9.} After this session I would trust the rover to do tasks I have not seen it do. {\footnotesize[Likert (1--5); id \texttt{arch\_l6}]}\par
\par\hspace{1em}{\footnotesize Rated 1--5: Strongly disagree, Disagree, Neutral, Agree, Strongly agree.}
\medskip\noindent\textbf{Q10.} Which moment changed your understanding the most, and what information would have helped you sooner? {\footnotesize[open-ended, optional; id \texttt{arch\_open2}]}\par
\medskip\noindent\textbf{Conf.} How confident are you in your answers in this section? {\footnotesize[confidence (1--5); id \texttt{arch\_conf}]}\par
\par\hspace{1em}{\footnotesize Rated 1--5: Guessing, Unsure, Somewhat, Fairly sure, Very sure.}
